\src{185}
One checks that the isomorphism (5.6) is precisely (5.2 bis) for $p=n$ after
setting $L=k$ by means of (5.7). Consequently, (5.2 bis) is an isomorphism for
$p=n$. To prove that it is an isomorphism for every $p$, it again suffices to
prove that, for $p<n$, the two sides of (5.2 bis) are coeffaceable functors of
$F$, and a fortiori (in view of Lemma 3) that both sides vanish when
$F=\Osh_X(-m)$ with $m$ large. For the first side this follows from Lemma 4; for
the second, using the duality theorem on $P$, we write
\[
  H^{p}(X,\Osh_X(-m))'
  =\operatorname{Ext}_{\Osh_P}^{r-p}
    (P;\Osh_X(-m),\Omega_P^{r}).
\]
The second side vanishes for $p<n$ and $m$ large, as follows from Lemma 5 (with
$X=P$ there) and from the fact that, as a coherent algebraic sheaf on $P$,
$\Osh_X$ has cohomological dimension $\leq r-n$ (since $X$ is nonsingular),
whence
\[
  \Extsh_{\Osh_P}^{r-p}(\Osh_X,\Omega_P^{r})=0
  \qquad\text{for }p<n.
\]

\fsection{6}{The duality theorem for singular varieties}

Let $X$ be a closed subset of dimension $n$ of projective space $P$ of
dimension $r$. Formula (5.6) must here be written
\begin{equation}\tag{6.1}
  H^{n}(X,F)'
  \simeq\operatorname{Hom}_{\Osh_X}(X;F,\omega_X^{n})
  =\operatorname{Ext}_{\Osh_X}^{0}(X;F,\omega_X^{n}),
\end{equation}
where we have set
\begin{equation}\tag{6.2}
  \omega_X^{n}=E^{0}(\Osh_X)
  =\Extsh_{\Osh_P}^{r-n}(\Osh_X,\Omega_P^{r}).
\end{equation}
As stated in Proposition 5, the sheaf so introduced does not in fact depend on
the chosen embedding of $X$ into a nonsingular variety $P$. Taking
$F=\omega_X^{n}$ in (6.1), we obtain
\begin{equation}\tag{6.2 bis}
  H^{n}(X,\omega_X^{n})'
  \simeq\operatorname{Hom}_{\Osh_X}(X;\omega_X^{n},\omega_X^{n}),
\end{equation}
\footnote{The printed text also numbers this formula ``(6.2)''; the suffix
``bis'' distinguishes the two occurrences.} and hence a distinguished element
of $H^{n}(X,\omega_X^{n})'$ corresponding to the identity homomorphism of
$\omega_X^{n}$:
\begin{equation}\tag{6.3}
  \gamma:H^{n}(X,\omega_X^{n})\longrightarrow k.
\end{equation}

